\chapter*{はじめに}
\addcontentsline{toc}{chapter}{はじめに}
\markboth{はじめに}{}
曲線で囲まれた図形の面積を求めたいとき，長方形や三角形の面積の公式だけでは困ることがある。
図形を細かく分ければ近い値は求められそうだが，どこまで分ければ正確になるのだろうか。
一方で，関数のグラフを描くときには，どこで増え，どこで減るかを知りたくなる。
この二つの問いを結び付けるのが，微分と積分である。

本書では，積分を中心に，高校で学ぶ極限・微分・積分の考え方を説明する。
面積を細かく分けるところから出発し，極限と微分を準備して，微分と積分がどのようにつながるのかを確かめる。
その後で，積分の計算と応用を考えていこう。

積分の問題を解く際には，公式を知っていても，何を使えばよいか分からないことがある。
例えば，平方根が含まれているというだけでは置換は一つに決まらない。
平方根の中の式を微分すると外側にある因子と一致するなら，その中身を新しい変数にするのが自然である。
円の方程式の形が見えるなら，三角関数で表す方法もある。
同じ「置換積分」であっても，簡単にしたい箇所が違うのである。
本書では，このように式を見て方針を決める理由を，実際の計算とともに示す。

また，一つの積分をいくつかの方法で解くと，異なる公式の関係が見えることがある。
三角関数の置換，円の幾何，双曲線関数は，別々に覚えるだけでは見えにくい関係をもっている。
基本の計算を身に付けた後では，その背景にも進めるようにした。
高度な結果を使わずに解く方法がある場合は，基本の道具でどこまで進めるかも確かめる。

\section*{対象読者と読む順序}
二次関数，整式の計算，三角関数の基本，指数・対数，数列の和を学んだ読者を想定する。
微積分を初めて学ぶ場合は，\chapterref{章:導入}から\chapterref{章:積分の応用}までを順に読み，章末問題で計算を確かめてほしい。
\chapterref{章:極限と微分}は後の説明に必要な微分をまとめているため，他の章より多くの内容を含む。
すでに微分を学んでいる場合は，\chapterref{章:定積分}から読み始め，必要になった箇所だけ\chapterref{章:極限と微分}へ戻ってよい。

\begin{center}
\begin{tabularx}{\linewidth}{@{}p{35mm}X@{}}
\toprule
読みたい内容 & 主な箇所\\
\midrule
積分の意味を知る & 第1〜3章。近似の和が定積分になり，原始関数で計算できる理由を扱う。\\
計算方法を選ぶ & 第4〜7章。式の整理，置換，部分積分，対称性を比較する。\\
積分を問題に使う & 第8〜9章。図形・運動・数列・不等式を通して，何を足すかを考える。\\
手法の背景へ進む & 第10〜11章。置換の幾何，逆関数，近似，広義積分を扱う。\\
自分で方針を決める & \chapterref{章:総合演習}。問題を解いてからヒントや巻末の解答を参照する。\\
\bottomrule
\end{tabularx}
\end{center}

\chapterref{章:積分と数列}の一般のベータ関数，第10〜11章，および\chapterref{章:積分の応用}の重心とパップス＝ギュルダンの定理は，高校の基本事項から先へ進むための発展内容である。
基本の学習に際して，これらの定理をすべて覚える必要はない。
発展内容では，高校の計算から何が広がるかと，新たに何を仮定するかを区別して読むとよい。
本書は高校数学全分野の教科書ではない。複素数平面，確率，整数などは，積分との関係が必要な範囲で扱う。

%
\section*{文字と説明について}
特に断らなければ，変数は実数を表す。
$\log x$は自然対数，$\e$はその底とし，三角関数の角は弧度法で測る。
$\mathbb{N}$を用いる場合は正の整数の集合とする。
不定積分に現れる$C$は積分定数であり，式ごとに値が変わってよい。
分母が零になる点などを除いて区間が分かれる場合には，それぞれの区間で定数を選ぶ。

定積分の$\dd x$は，どの変数について積分しているかを示す。
置換の途中で$\dd u=g'(x)\dd x$と書くときには，置換積分の公式を簡潔に表している。
この記号だけを根拠に，どのような操作でも分数のように行えると考えないようにしたい。

重要な公式は，その公式を使う理由と導出を本文に示した。
一方，連続関数の積分可能性など，実数の性質にまで戻る証明は，採用する事実として明示している。
発想を得るための近似と，結果が正しいことを確かめる議論も分けてある。
例題と章末問題は本書の説明のために設定したものであり，特定の入試問題の転載ではない。
章末問題と総合演習には，巻末にすべて解答を付けた。
%
